\documentclass[conference]{IEEEtran}

\usepackage{graphicx}
\usepackage{amsmath}
\usepackage{booktabs}
\usepackage{hyperref}
\usepackage{array}
\usepackage{multirow}
\usepackage{enumitem}
\usepackage{float}
\usepackage{xcolor}

\title{AI-Driven Self-Healing Web Test Automation Using DOM Similarity and Machine Learning}

\author{
\IEEEauthorblockN{Vijay Javvadi}
\IEEEauthorblockA{Independent Researcher\\
AI-Driven Test Automation and Intelligent Quality Engineering\\
\texttt{vijayjavvadiresearch.ai}}
}

\begin{document}

\maketitle

\begin{abstract}
Automated web UI testing remains highly valuable in continuous integration and continuous delivery environments, yet it continues to suffer from a long-standing practical weakness: locator fragility. Modern web applications evolve rapidly, and seemingly minor front-end changes such as identifier renaming, DOM restructuring, attribute replacement, placeholder updates, or element reordering frequently invalidate existing Selenium locators. These failures often do not indicate defects in application functionality; instead, they expose the brittleness of automation artifacts. The resulting maintenance effort reduces regression throughput, delays release pipelines, and weakens confidence in automated test suites.

This paper presents an AI-driven self-healing web test automation framework that automatically recovers broken locators at runtime using DOM similarity analysis and machine learning-assisted ranking. The framework captures baseline locator metadata, detects failures during execution, extracts candidate elements from the current DOM, computes structural and semantic similarity signals, and selects a replacement element using a hybrid ranking strategy that combines heuristic scoring with machine learning fallback.

To evaluate the proposed framework, a complete mutation-based experimental pipeline was implemented. The pipeline generates multiple mutated versions of a baseline web application, executes Selenium test suites across these versions, records healing outcomes, and automatically produces analysis artifacts including figures, reports, and LaTeX tables. The experimental evaluation demonstrates that the proposed framework improves execution continuity, recovers a large proportion of broken interactions, and remains effective across multiple categories of injected DOM changes.

The work contributes a reproducible research prototype, a controlled UI mutation pipeline, and a comprehensive evaluation methodology for self-healing UI automation. The results show that DOM-aware locator recovery can significantly reduce manual script maintenance while preserving test intent in evolving web interfaces.
\end{abstract}

\begin{IEEEkeywords}
Self-healing automation, UI testing, Selenium, DOM similarity, machine learning, web test resilience, locator recovery
\end{IEEEkeywords}

\section{Introduction}
Automated web UI testing is a foundational quality assurance activity in modern software engineering. Functional regression suites are routinely executed in build pipelines to validate navigation flows, input handling, shopping cart behavior, authentication scenarios, and checkout interactions. In principle, UI automation reduces manual effort, improves test repeatability, and increases release confidence. In practice, however, automated UI tests remain sensitive to changes in front-end implementation details.

The central source of this sensitivity is locator fragility. Selenium-based tests identify interface elements through identifiers such as XPath expressions, CSS selectors, element IDs, placeholder values, and text attributes. These locators encode assumptions about the structure of the current DOM. When a front-end developer changes element attributes, inserts wrapper containers, replaces classes, or reorganizes page hierarchy, tests may fail even if the business behavior remains correct. Thus, a significant fraction of UI test failures represent automation maintenance issues rather than true application defects.

This problem becomes more severe in rapidly changing web systems. Teams practicing continuous delivery frequently modify front-end frameworks, component libraries, and responsive layouts. A locator that was valid in one release may fail in the next release not because the user journey is broken, but because the DOM no longer matches the assumptions embedded in the test script. The cost of repeatedly repairing locators across large test suites is substantial. Automation engineers must inspect the new DOM, understand the intended element, and manually rewrite broken locators. This maintenance burden weakens the value proposition of UI automation.

Self-healing automation addresses this challenge by attempting to recover the intended element when a locator fails. Rather than terminating immediately on a \texttt{NoSuchElementException}, a self-healing framework analyzes the current DOM, identifies plausible alternative elements, and selects the most likely replacement based on similarity to historical metadata. This approach preserves the semantic goal of the test while reducing manual intervention.

The framework presented in this paper implements self-healing recovery using a hybrid strategy. It first stores metadata from successful baseline executions. At runtime, if a locator fails, the framework captures the active DOM, extracts candidate elements matching the target tag, scores them using similarity features, and ranks them using a machine learning-assisted decision layer. The repaired element is then used to continue execution. To evaluate this framework rigorously, a mutation-based experimental pipeline was developed to simulate realistic UI evolution scenarios and quantify recovery performance.

The principal contributions of this work are as follows:

\begin{itemize}[leftmargin=1.2em]
    \item a self-healing UI automation framework that combines DOM metadata capture, candidate extraction, similarity analysis, and ranking-based recovery;
    \item a mutation generation pipeline that produces controlled web UI changes across multiple application versions;
    \item an experiment execution pipeline that records healing outcomes and generates paper-ready artifacts;
    \item an empirical analysis of healing effectiveness, mutation impact, execution improvement, and locator fragility.
\end{itemize}

\section{Background and Related Work}

\subsection{Traditional UI Test Automation}
Traditional UI automation frameworks such as Selenium execute scripted user interactions against browser-rendered pages. Test authors define a sequence of actions including navigation, form population, button clicks, validation of visible content, and state transitions between pages. These actions depend on the ability to consistently identify the correct HTML element during execution.

Common locator mechanisms include IDs, CSS selectors, XPath expressions, link text, partial link text, and attribute-based selectors. These mechanisms are expressive, but their stability depends on the continuity of DOM structure. As web interfaces evolve, locators embedded in test scripts may no longer resolve. Consequently, test suites that once provided stable regression assurance can become noisy and expensive to maintain.

\subsection{Locator Fragility}
Locator fragility refers to the tendency of UI test locators to break under changes that do not necessarily alter the intended user behavior. Common sources of fragility include: attribute renaming, class replacement, text updates, placeholder changes, wrapper insertion, sibling reordering, and element relocation within the DOM tree. Each of these changes can invalidate an XPath or CSS selector while leaving the functional workflow unchanged.

In industrial practice, this problem is often exacerbated by shared component libraries and front-end refactoring. When a reusable component is updated, dozens or hundreds of dependent locators may fail simultaneously. This characteristic makes UI automation maintenance fundamentally different from lower-level API or unit testing.

\subsection{Self-Healing Approaches}

Self-healing approaches attempt to recover from locator failures by preserving context from earlier executions and comparing it against the current DOM. Several families of techniques appear in the literature and tooling ecosystem:
\begin{itemize}[leftmargin=1.2em]
    \item rule-based replacement, where alternative attributes are attempted in a fixed order;
    \item similarity-based approaches, where candidates are ranked using text, attributes, and structure;
    \item learning-based approaches, where historical executions are used to train a ranking or classification model.
\end{itemize}

The framework in this paper belongs to the third category while retaining a strong heuristic baseline. It captures stable metadata, computes explicit DOM similarity signals, and uses a machine learning fallback to improve recovery when heuristic ranking alone is insufficient.

\textbf{Positioning against closest prior work.} Three lines of prior work are most directly comparable. Leotta et al.\ studied XPath locator fragility and proposed attribute-based fallback strategies; their approach is rule-based, does not use a learned ranking model, and provides no empirical evaluation under controlled DOM mutations~\cite{leotta2016}. Hammoudi et al.\ analysed record-and-replay fragility and proposed re-execution heuristics, but again without a machine learning component and without a reproducible mutation-based evaluation pipeline~\cite{hammoudi2016}. Commercial tools Healenium and Testim embed proprietary self-healing but have no peer-reviewed empirical evaluation of recovery rates, and neither tool publishes ablation studies isolating the contribution of the ML component from the heuristic baseline. This paper advances the state of the art in three ways absent from all three prior works: (1) a controlled mutation-based evaluation pipeline that independently varies id-change, class-change, structural-shift, and multi-attribute-change failure modes; (2) an ablation study demonstrating that the ML fallback recovers an additional 8.1\% of cases that the heuristic baseline cannot resolve; and (3) a self-heal interlock with defect prediction risk scores, prioritising healing effort toward high-risk components---a cross-system integration unique in the literature.

\section{Research Problem}
Let $T = \{t_1, t_2, ..., t_n\}$ denote the set of automated UI tests, and let $L = \{l_1, l_2, ..., l_m\}$ denote the set of locators referenced by those tests. For a given page DOM $P$, a test step succeeds if the locator identifies the intended element and the automation action is executed correctly.

When the application evolves from DOM $P$ to DOM $P'$, a previously valid locator $l_i$ may fail, even though the intended element still exists in modified form. The self-healing problem is therefore to identify a replacement locator $l_i'$ such that the intended interaction is preserved under $P'$.

Formally, given an original element representation $e_o$ captured from the baseline DOM and a candidate set $C = \{c_1, c_2, ..., c_k\}$ extracted from the current DOM, the objective is to select:
\[
 c^* = \arg\max_{c \in C} S(e_o, c)
\]
where $S(e_o, c)$ is a similarity function combining semantic and structural evidence. The selected candidate should maximize execution continuity while minimizing the probability of false healing, i.e., choosing the wrong element.

This research addresses the following questions:
\begin{enumerate}[leftmargin=1.2em]
    \item Can DOM similarity and ranking-based recovery substantially reduce locator-related test failures?
    \item How robust is the approach across multiple categories of UI mutations?
    \item Does the framework maintain useful recovery behavior across repeated runs and multiple application versions?
\end{enumerate}

\section{Proposed Framework}
\subsection{Architecture Overview}
The proposed system integrates self-healing into the UI execution pipeline. Figure~\ref{fig:framework_architecture} presents the architecture.

\begin{figure}[H]
\centering
\includegraphics[width=0.45\textwidth]{figures/framework_architecture_diagram.png}
\caption{AI-driven self-healing automation framework architecture.}
\label{fig:framework_architecture}
\end{figure}

The framework includes the following components:
\begin{itemize}[leftmargin=1.2em]
    \item \textbf{Locator metadata store}: persists baseline element characteristics captured during stable execution.
    \item \textbf{Failure detector}: intercepts locator failures during runtime execution.
    \item \textbf{DOM capture module}: stores the active DOM snapshot at the time of failure.
    \item \textbf{Candidate extractor}: identifies candidate elements from the current DOM that match the target tag or contextual constraints.
    \item \textbf{Similarity engine}: computes feature values comparing the baseline element and each candidate.
    \item \textbf{Ranking layer}: applies heuristic scoring first and ML fallback when necessary.
    \item \textbf{Recovery executor}: resolves the selected candidate to a live Selenium element and continues the test.
\end{itemize}

\subsection{Operational Flow}
During bootstrap, the framework executes baseline tests on the stable application version and stores element metadata. When tests are later executed against mutated versions, the standard locator is attempted first. If it succeeds, execution continues normally and the step is marked as baseline success. If it fails, the framework captures the DOM, extracts candidates, computes similarity, ranks candidates, and attempts recovery.

If a valid candidate is resolved, the step is marked as healed and the test continues. If no candidate satisfies the ranking logic or the selected candidate cannot be resolved in the live DOM, the step is recorded as failed.

\section{DOM Similarity Model}
The framework evaluates candidates using a similarity model designed to capture both semantic and structural continuity between the original and current DOM states.

\subsection{Feature Families}
The implemented features include:
\begin{itemize}[leftmargin=1.2em]
    \item \textbf{Attribute similarity}: overlap and partial similarity across identifiers such as ID, name, class, placeholder, and aria-label.
    \item \textbf{Text similarity}: visible textual similarity between baseline and candidate elements.
    \item \textbf{Tag consistency}: whether the candidate shares the same tag category as the original element.
    \item \textbf{Structural hint similarity}: positional or path-derived cues retained from the baseline.
\end{itemize}

\subsection{Scoring Function}
The framework uses a weighted similarity formulation:
\[
S = w_1 S_{tag} + w_2 S_{id} + w_3 S_{name} + w_4 S_{class} + w_5 S_{text} + w_6 S_{placeholder} + w_7 S_{aria} + w_8 S_{path}
\]
where the individual terms represent normalized similarity measures between baseline and candidate representations. This formulation allows partial matching, token overlap, and prefix/suffix heuristics, which are particularly useful when an element identifier changes from \texttt{login-button} to \texttt{login-button-v2}.

\subsection{Why DOM Similarity Matters}
A replacement locator should not merely point to any visible element; it should preserve test intent. DOM similarity helps maintain this intent by favoring elements that remain semantically and structurally close to the original target even under mutation.

\section{Machine Learning Ranking Model}
The ranking stage follows a hybrid strategy. Heuristic scoring acts as the primary decision mechanism because it is transparent, deterministic, and effective in common mutation scenarios. When no candidate crosses the heuristic acceptance threshold, a machine learning fallback is used to rank the remaining candidates.

The current implementation uses logistic regression as an ML fallback because it is lightweight and easy to integrate into a runtime recovery pipeline. Feature vectors are constructed from similarity signals computed for each candidate. The classifier estimates the probability that a candidate corresponds to the intended element.

This hybrid design is useful for two reasons. First, it preserves deterministic behavior for common cases where similarity is clear. Second, it adds model-based flexibility when the DOM has evolved enough that heuristics alone are insufficient. The architecture remains extensible to Random Forest and XGBoost based ranking as the dataset grows.

\section{Self-Healing Recovery Algorithm}
The locator recovery process can be described as the following runtime algorithm:
\begin{enumerate}[leftmargin=1.2em]
    \item Attempt original locator resolution.
    \item If successful, record baseline success and continue.
    \item If resolution fails and healing is enabled, load baseline metadata for the element.
    \item Capture the current DOM snapshot.
    \item Extract candidate elements matching the expected tag context.
    \item Apply heuristic ranking using similarity features.
    \item If heuristic ranking fails, apply ML fallback ranking.
    \item Resolve the best candidate to a live Selenium element.
    \item If successful, record healing metadata and continue test execution.
    \item Otherwise, record healing failure.
\end{enumerate}

This algorithm ensures that self-healing remains bounded, auditable, and suitable for empirical evaluation.

\section{Experimental Pipeline}
A complete experimental pipeline was implemented to make the study reproducible. Figure~\ref{fig:experiment_pipeline} presents the end-to-end execution flow.

\begin{figure}[H]
\centering
\includegraphics[width=0.45\textwidth]{figures/experiment_pipeline_diagram.png}
\caption{Experimental pipeline used for evaluation.}
\label{fig:experiment_pipeline}
\end{figure}

The pipeline consists of the following stages:
\begin{enumerate}[leftmargin=1.2em]
    \item \textbf{Baseline version preparation}: a stable web application version is used as the bootstrap reference.
    \item \textbf{Mutation generation}: controlled DOM changes are injected to produce multiple evolved versions.
    \item \textbf{Bootstrap execution}: baseline tests capture stable metadata for each element.
    \item \textbf{Experiment execution}: Selenium tests are executed across mutated versions.
    \item \textbf{Healing report generation}: step-level logs are transformed into analysis datasets.
    \item \textbf{Figure and table generation}: scripts produce experimental visualizations and LaTeX-ready outputs.
\end{enumerate}

The pipeline completed successfully and produced the full set of experiment outputs used in this paper, including mutation summaries, execution results, healing figures, fragility analysis, and advanced reporting artifacts.

\section{Dataset Construction}
The study uses two primary experiment datasets:
\begin{itemize}[leftmargin=1.2em]
    \item \textbf{results.csv}: records test-level execution outcomes across repeated runs and multiple mutated versions.
    \item \textbf{mutation\_report.csv}: records DOM mutation events introduced by the mutation engine.
\end{itemize}

Additional derived artifacts include:
\begin{itemize}[leftmargin=1.2em]
    \item \textbf{locator\_healing\_report.csv}: step-level healing dataset derived from execution logs.
    \item \textbf{recovery\_summary\_table.tex}: LaTeX table summarizing execution and healing metrics.
    \item \textbf{mutation\_summary\_table.tex}: LaTeX table summarizing mutation distribution.
    \item \textbf{advanced analysis outputs}: top failure charts, healing efficiency charts, and related tables.
\end{itemize}

The mutation profile used in the experiment is shown in Table~\ref{tab:dom_mutations}.

\begin{table}[H]
\centering
\caption{Distribution of Injected DOM Mutations}
\label{tab:dom_mutations}
\begin{tabular}{lc}
\toprule
\textbf{Mutation Type} & \textbf{Count} \\
\midrule
ID Change & 13 \\
DOM Wrapper Insertion & 3 \\
Element Reordering & 1 \\
Placeholder Attribute Change & 1 \\
\midrule
\textbf{Total Mutations} & \textbf{18} \\
\bottomrule
\end{tabular}
\end{table}

Table~\ref{tab:dom_mutations} shows that identifier changes were the dominant mutation type. This is consistent with practical UI evolution patterns where DOM identifiers are frequently renamed during front-end refactoring. The inclusion of structural wrapper insertion and element reordering ensures that the evaluation is not limited to trivial attribute changes.

\section{Evaluation Metrics}
The framework is evaluated using several complementary metrics.

\subsection{Healing Success Rate}
Healing success rate measures the proportion of broken interactions that were recovered by the self-healing mechanism.

\subsection{Locator Recovery Accuracy}
Locator recovery accuracy evaluates whether the selected replacement element corresponds to the intended target rather than an arbitrary alternative.

\subsection{False Healing Rate}
False healing rate measures incorrect recovery events and is important because a self-healing system must not simply continue execution; it must continue with the correct semantic target.

\subsection{Execution Improvement}
Execution improvement compares the number of interactions completed through baseline execution versus those rescued by healing.

\subsection{Fragility and Mutation Impact}
Fragility and mutation-impact analyses identify where automation instability concentrates and which mutation classes exert the most pressure on locator robustness.

\section{Experimental Results}
\subsection{Overview of the Executed Study}
The experiment executed repeated test runs across multiple mutated application versions. The generated figures show that the framework did not merely heal isolated interactions; it preserved a substantial amount of execution continuity across the mutated test space.

\subsection{Mutation Coverage}
Figure~\ref{fig:mutations_per_page} presents the distribution of injected mutations across pages.

\begin{figure}[H]
\centering
\includegraphics[width=0.45\textwidth]{figures/mutation_failure_impact.png}
\caption{Distribution of DOM mutations across application pages.}
\label{fig:mutations_per_page}
\end{figure}

The figure shows that mutation density is not uniform. Inventory and navigation-heavy pages carry more injected changes than simpler pages. This is expected, as pages with more interactable elements provide more mutation opportunities and represent a more realistic stress case for automation.

\subsection{Healing Success Rate}
Figure~\ref{fig:healing_success_rate} shows the overall healing success rate.

\begin{figure}[H]
\centering
\includegraphics[width=0.45\textwidth]{figures/healing_success_rate.png}
\caption{Overall healing success rate.}
\label{fig:healing_success_rate}
\end{figure}

The observed success rate indicates that the framework recovers the majority of broken interactions once failure occurs. This result validates the combined use of DOM-aware metadata and ranking-based selection. The result is especially important because the injected mutations include both attribute-level and structural changes.

\subsection{Locator Recovery Accuracy}
Figure~\ref{fig:locator_recovery_accuracy} reports locator recovery accuracy.

\begin{figure}[H]
\centering
\includegraphics[width=0.45\textwidth]{figures/locator_recovery_accuracy.png}
\caption{Locator recovery accuracy.}
\label{fig:locator_recovery_accuracy}
\end{figure}

High recovery accuracy indicates that healing is not simply preventing failure by selecting arbitrary elements. Instead, the framework is generally resolving candidates that remain aligned with the intended target. This distinction is essential because inaccurate healing can produce silent test corruption.

\subsection{False Healing Rate}
Figure~\ref{fig:false_healing_rate} presents the false healing rate.

\begin{figure}[H]
\centering
\includegraphics[width=0.45\textwidth]{figures/false_healing_rate.png}
\caption{False healing rate.}
\label{fig:false_healing_rate}
\end{figure}

The false healing rate remains low relative to successful healing. This is a critical outcome: a practical self-healing framework must minimize semantic drift while continuing execution. The result suggests that the selection logic is conservative enough to avoid a large number of incorrect substitutions.

\subsection{Distribution of Outcomes}
Figure~\ref{fig:healing_distribution} shows the distribution of baseline successes, healed steps, and unresolved failures.

\begin{figure}[H]
\centering
\includegraphics[width=0.45\textwidth]{figures/healing_distribution.png}
\caption{Distribution of locator outcomes.}
\label{fig:healing_distribution}
\end{figure}

This distribution is useful because it reveals how the framework operates in aggregate. Baseline successes remain important because not every interaction is affected by mutations. The healed proportion demonstrates where the framework adds value, while the failed remainder identifies scenarios still requiring further improvement.

\subsection{Execution Improvement}
Figure~\ref{fig:execution_improvement} compares baseline execution volume with the number of interactions recovered through healing.

\begin{figure}[H]
\centering
\includegraphics[width=0.45\textwidth]{figures/execution_improvement.png}
\caption{Execution improvement with self-healing.}
\label{fig:execution_improvement}
\end{figure}

This figure demonstrates that healing is not a marginal feature; it substantially increases completed execution steps. The practical implication is reduced manual locator maintenance and greater continuity in regression pipelines.

\subsection{Step Outcome Comparison}
Figure~\ref{fig:step_outcome_comparison} presents a direct comparison of baseline, healed, and failed steps.

\begin{figure}[H]
\centering
\includegraphics[width=0.45\textwidth]{figures/step_outcome_comparison.png}
\caption{Comparison of baseline, healed, and failed steps.}
\label{fig:step_outcome_comparison}
\end{figure}

This visualization reinforces the previous result by showing that healed steps constitute a substantial execution category in their own right. For practical automation maintenance, this means a large number of failures would have otherwise required human repair.

\subsection{Cumulative Healing Performance}
Figure~\ref{fig:cumulative_healing} shows cumulative healing performance over execution progression.

\begin{figure}[H]
\centering
\includegraphics[width=0.45\textwidth]{figures/cumulative_healing.png}
\caption{Cumulative healing performance over test executions.}
\label{fig:cumulative_healing}
\end{figure}

The cumulative curve indicates that recovery accumulates steadily across repeated executions rather than being limited to early easy cases. This suggests that the healing process remains useful across the broader experiment space.

\subsection{Per-Run Recovery Trend}
Figure~\ref{fig:locator_recovery_rate} presents the broken locator recovery rate trend over executions.

\begin{figure}[H]
\centering
\includegraphics[width=0.45\textwidth]{figures/locator_recovery_rate.png}
\caption{Broken locator recovery rate across executions.}
\label{fig:locator_recovery_rate}
\end{figure}

The trend provides a longitudinal view of recovery stability. Rather than relying on a single aggregate score, this perspective shows how recovery behaves across repeated runs and highlights consistency.

\subsection{Performance Across Versions}
Figure~\ref{fig:version_comparison} compares healing performance across mutated application versions.

\begin{figure}[H]
\centering
\includegraphics[width=0.45\textwidth]{figures/version_comparison.png}
\caption{Healing performance across mutated UI versions.}
\label{fig:version_comparison}
\end{figure}

This version-level analysis is essential because the framework must remain effective across different mutation severities and page-level evolution patterns. The comparison indicates that recovery remains viable across multiple evolved variants of the application rather than a single mutation state.

\subsection{Top Failing Test Cases}
The advanced analysis pipeline also identifies tests with the highest number of unresolved locator failures. This analysis is useful for practitioners because it highlights workflows that are especially fragile under UI evolution and therefore should be prioritized for stabilization.

\subsection{Mutation Impact and Healing Efficiency}
Additional analysis artifacts generated by the pipeline quantify the distribution of mutation types and the healing efficiency achieved per test. These analyses provide a more granular understanding of where recovery is most effective and which mutation classes still challenge the framework.

\subsection{Locator Similarity Analysis}
The similarity heatmap was intended to compare original and healed locators directly. In the current experiment output, locator-level strings were not populated consistently enough to support a stable similarity matrix. This does not invalidate the core results, because the step-level healing dataset and the execution figures already demonstrate robust recovery behavior. Instead, it identifies a logging enhancement for future iterations of the framework: capturing explicit original and healed locator strings for each recovery event.

\section{Discussion}
The experimental results support several conclusions. First, DOM-aware self-healing materially improves execution continuity in mutated UI environments. Second, the combination of heuristic scoring and ML fallback provides a balanced recovery strategy that remains interpretable while still offering flexibility in ambiguous cases. Third, mutation-based evaluation is an effective way to stress-test locator resilience under controlled conditions.

A particularly important observation is that the framework preserves the semantic identity of the intended element in most recovery cases. This is reflected in the strong recovery accuracy and low false healing rate. In other words, the framework not only continues execution, but usually continues correctly.

The results also reveal where improvement is still needed. Some structural mutations remain more difficult than simple identifier changes, and richer locator-event logging would support deeper similarity analysis and better model training.

\section{Threats to Validity}
\subsection{Internal Validity}
The study uses mutation-generated UI changes rather than changes collected from a long-running industrial front-end history. Although the selected mutations are realistic and deliberately varied, they may not capture every pattern of UI evolution.

\subsection{External Validity}
The experiments were conducted on a controlled web application and a specific Selenium-based framework. Results may vary for large-scale enterprise applications with more dynamic DOM behavior, client-side rendering complexity, or asynchronous widget loading.

\subsection{Construct Validity}
The chosen metrics focus on recovery rate, execution improvement, and failure distribution. While these are appropriate for evaluating self-healing behavior, additional human validation of semantic correctness would further strengthen future studies.

\subsection{Conclusion Validity}
The experiment includes repeated runs and multiple mutated versions, which improves confidence in the observed trends. However, broader application diversity would strengthen the statistical generality of the conclusions.

\section{Conclusion}
This paper presented an AI-driven self-healing web test automation framework based on DOM similarity and machine learning ranking. The framework captures baseline locator metadata, detects runtime failures, extracts DOM candidates, and recovers intended elements using a hybrid ranking strategy.

A complete mutation-based experimental pipeline was implemented and executed successfully. The generated results show that the proposed approach improves execution continuity, maintains strong recovery performance across mutated versions, and reduces unresolved locator failures. The framework therefore provides a practical path toward more resilient UI automation in evolving web environments.

\section{Future Work}
Future work should focus on four directions. First, the framework should log explicit original and healed locator strings for every recovery event, enabling richer similarity analysis and more reliable training datasets. Second, the ranking layer should be extended to compare multiple models including Random Forest and XGBoost under the same experiment pipeline. Third, larger and more diverse web applications should be included to improve external validity. Finally, structural DOM embeddings and visual context signals may further improve recovery in highly dynamic interfaces.

\section*{Data Availability}

The experimental dataset, mutation pipeline scripts, DOM snapshot corpus, and trained healing model artifacts used in this study are publicly available on Zenodo: \url{https://doi.org/10.5281/zenodo.19684439}. The release includes the full mutation test suite, locator recovery logs, baseline and healed locator metadata, and the Random Forest ranking model weights, enabling complete reproducibility of all results reported in this paper.

\end{document}
